\section{Native normalization, dimension and cover criteria}
\label{sec:di-reused-native-decomposition-criteria}

The following criteria are applied to the original coefficient maps,
form sheaves, finite sources and restrictions throughout the decomposition.

\begin{lemma}[Normalization over the original coefficient field]
\label{lem:di-reused-native-normalization}
\lean{EnvelopeProperty.Cartier.DeligneIllusieAlternatingNormalization.normalizationScalar,EnvelopeProperty.Cartier.DeligneIllusieAlternatingNormalization.factorial_cast_ne_zero,EnvelopeProperty.Cartier.DeligneIllusieAlternatingNormalization.normalizationScalar_mul_factorial,EnvelopeProperty.Cartier.DeligneIllusieAlternatingNormalization.normalizedAlternatization,EnvelopeProperty.Cartier.DeligneIllusieAlternatingNormalization.normalizedAlternatization_apply,EnvelopeProperty.Cartier.DeligneIllusieAlternatingNormalization.comp_normalizedAlternatization,EnvelopeProperty.Cartier.DeligneIllusieAlternatingNormalization.normalizedAlternatization_closed,EnvelopeProperty.Cartier.DeligneIllusieAlternatingNormalization.normalizedAlternatization_recovery,EnvelopeProperty.Cartier.DeligneIllusieAlternatingNormalization.normalizedAlternatization_naturality}
\leanok
Let $k$ be a field of characteristic $p$, and let $\beta:k\to A$ be
the original structural homomorphism to a commutative ring. The scalar
$\beta((q!)^{-1})$ normalizes the signed sum over permutations of a
$q$-multilinear map. The resulting map is alternating, commutes with
postcomposition and compatible semilinear restriction, and remains
closed when the original values are closed. If $q<p$, its value on an
already alternating map is that original map.
\end{lemma}
\begin{proof}
\leanok
Each signed permutation is multilinear, and the antisymmetrization
universal property gives the alternating map. Postcomposition and
compatible restriction distribute over the finite sum and preserve
the original structural scalar. An outgoing linear differential also
distributes over it, giving closedness. If the original map is
alternating, its value on a permuted tuple contributes the permutation
sign a second time, so every signed summand has the same value. There
are $q!$ summands. For $q<p$ none of the factors $1,\ldots,q$ vanishes
in $k$, so $q!\ne0$ and its inverse cancels it in $k$. Applying the
original $\beta$ gives the cancellation in $A$, including when $A$ is
the zero ring. No exact-characteristic or nonemptiness assumption on
$A$ is needed.
\end{proof}

\begin{lemma}[Original geometric dimension support]
\label{lem:di-reused-native-dimension-support}
\lean{EnvelopeProperty.Cartier.DeligneIllusieDimensionBound.exteriorCotangentSheaf_isZero,EnvelopeProperty.Cartier.DeligneIllusieDimensionBound.twistStructure,EnvelopeProperty.Cartier.DeligneIllusieDimensionBound.twistStructure_smooth,EnvelopeProperty.Cartier.DeligneIllusieDimensionBound.twist_exteriorCotangentSheaf_isZero,EnvelopeProperty.Cartier.DeligneIllusieDimensionBound.formModule_isZero,EnvelopeProperty.Cartier.DeligneIllusieDimensionBound.complex_extend_isStrictlyLE,EnvelopeProperty.Cartier.DeligneIllusieDimensionBound.goodTruncationInclusion,EnvelopeProperty.Cartier.DeligneIllusieDimensionBound.goodTruncationInclusion_isIso,EnvelopeProperty.Cartier.DeligneIllusieDimensionBound.goodTruncationInclusion_quasiIso}
\leanok
Let $f:X\to\operatorname{Spec}k$ be smooth of pure relative dimension
$n$. Original form sheaves vanish in degrees above $n$. The true
Frobenius base change has the same dimension and the same vanishing.
The original pushed de Rham complex, extended by zero, is strictly
supported in degrees at most $n$. If $n<p$, its original good-truncation
inclusion at $p-1$ is an isomorphism, hence a quasi-isomorphism.
\end{lemma}
\begin{proof}
\leanok
On every standard smooth affine chart, the actual cotangent module is
free with its original $n$-element basis. Its $q$th exterior power has
the basis indexed by $q$-element subsets; there are none for $q>n$.
The original affine exterior-sheaf comparison and open-cover detection
give the sheaf vanishing. Smooth relative dimension is stable under
the actual Frobenius base change, so the same argument applies to the
true twist. Pushforward sends zero modules to zero modules. Extension
by zero adds zero terms in negative degrees. Finally $n<p$ places all
nonzero terms below the cutoff, including a zero outgoing differential
at its top. The original good-truncation inclusion is therefore an
isomorphism in every degree.
\end{proof}

\begin{lemma}[The original bounded-source inclusion]
\label{lem:di-reused-native-dimension-source}
\lean{EnvelopeProperty.Cartier.DeligneIllusieDimensionSource.source,EnvelopeProperty.Cartier.DeligneIllusieDimensionSource.sourceInclusion,EnvelopeProperty.Cartier.DeligneIllusieDimensionSource.sourceInclusion_isIso,EnvelopeProperty.Cartier.DeligneIllusieDimensionSource.sourceIso,EnvelopeProperty.Cartier.DeligneIllusieDimensionSource.ι_sourceIso_hom}
\leanok
\uses{lem:di-reused-native-dimension-support}
For the original true twist put
$S_a=\bigoplus_{q<a}\Omega^q[-q]$ in natural degrees. If $n<p$, the
original coproduct inclusion $S_{n+1}\to S_p$ is an isomorphism. Its
forward map sends each original smaller summand to the original larger
summand in exactly the same degree.
\end{lemma}
\begin{proof}
\leanok
Define its inverse by projecting each summand of degree at most $n$
to the corresponding smaller summand and sending the remaining
summands to zero. The original form sheaves in these remaining degrees
are zero by the dimension criterion. Each composite is consequently
the identity on every original summand. The coproduct universal
property proves both inverse laws and the stated inclusion formula.
\end{proof}

\begin{lemma}[Detection by the original cover]
\label{lem:di-reused-native-cover-detection}
\lean{EnvelopeProperty.Cartier.DeligneIllusieCechLocalQuasiIso.isIso_of_restrictions,EnvelopeProperty.Cartier.DeligneIllusieCechLocalQuasiIso.quasiIsoAt_of_restrictions,EnvelopeProperty.Cartier.DeligneIllusieCechLocalQuasiIso.quasiIso_of_restrictions}
\leanok
An original module-sheaf map is invertible if its restrictions to an
actual open cover are invertible. A map of the original integer
cochain complexes is a quasi-isomorphism in degree $q$, or in all
degrees, if the same property holds after every original cover
restriction.
\end{lemma}
\begin{proof}
\leanok
At every point select a cover member containing it. The actual open
restriction identifies its stalk with the original stalk, so the
restricted inverse gives an inverse on that stalk. Stalks detect
isomorphisms of the underlying sheaves, and the module-sheaf forgetful
functor reflects isomorphisms. Open restriction preserves the original
homology. Its native homology naturality square therefore identifies
the restricted original homology map with the homology map of the
restricted complex. Apply the sheaf-map criterion in degree $q$, and
then in every degree for the whole quasi-isomorphism criterion.
\end{proof}
